\documentclass[a4paper]{article}

\title{Sonatas}
\author{Sadie}

\newcommand\fl{$\flat$}
\newcommand\sh{$\sharp$}

\begin{document}
\maketitle

This is a suggested sequel to \textit{24 Preludes and Fugues}
and is offered for a small number of ambitious writers willing to
write a longer piece.

In my humble opinion, \textit{Sonata Form} is the greatest musical form
there is, and deserves to be better known generally (i.e.~outside
music).  It was developed early in the classical period (the time of
Haydn and Mozart) and is still used today.  The first movement of a
symphony is always in sonata form, and the last movement often is too.
String quartets and piano sonatas usually use sonata form too.
Overtures are often in sonata form.

This extended prompt will take you through the ideas and structure of
a somewhat generalised sonata form.  No musical knowledge is required.
Your task is to write a piece, a long short-story, a short novella, or
the equivalent length as poetry (a chap book perhaps) or a play or
piece of journalism, in sonata form.

\section{The idea}

We start with keys and subject, as for the preludes and fugues.
From a literary perspective, a key is a short list of adjectives
that describe a character trait, mood or emotion.  And the themes
are objects (e.g.~people, places, things, even abstract things)
that represent this key or have this key.  You will see below that
single themes may exist in different keys, as for example
a happy child, a sad child, and an excited child.

Sonata form has two main keys, called tonic and dominant (think
of these as `home' and `away')
and two main themes, called \textit{first subject} and
\textit{second subject}, one in each of these keys.  The idea
is that these  are in conflict.  There are three main sections,
\textit{exposition} where the themes and this conflict are
introduced, \textit{development} where the themes and conflict
are explored in detail, often passing through many other keys
as well as the home and away keys, and \textit{recapitulation}
where the first subject in the home key is repeated, followed by the
second theme but this time \textit{also in the home key} as a way
to indicate that the conflict has been resolved.

Thus sonata form is about setting up a conflict between two
themes and keys, exploring that conflict, and then resolving it.
Traditionally, the first subject is often fiery and the second
lyrical and smooth, but there is no need to conform to this.

\clearpage

\section{The structure}

As just stated, there are two main themes and three main sections to a
piece in sonata form, and some options.  It is possible to introduce
additional sections (introduction and coda) and a third theme
(transition).  However, the first and second subjects are the
main themes and the main focus.

Here is a mapping of an entire piece in sonata form
with the three main sections, optional sections
and the subsections within them.

\begin{itemize}
\item (Optional) Introduction.  This in typically short and may but
  need not be related to the main subjects.  In music, an introduction
  if present is usually in a slow tempo.
\item Exposition. Usually divided into three subsections.
\begin{itemize}
\item First subject.  A presentation of the first subject in the
  home key.  This can be a straightforward account of it.  For more
  advanced and later pieces (late 19thC and early 20thC) the
  first subject is a group of themes.  Whatever is chosen, it
  should be all recognisable as `first subject material'.
\item Transition.  As it sounds, this is material that leads out
  of the first subject and into the second subject.  It is usually
  quite short and can be quite simple, or it might have associated
  with it a theme of its own.  But primarily it is a period
  of transition, as its name suggests.
\item Second subject.  A presentation of the second subject in the
  away key.  Again, this can be a straightforward account
  or a group of themes.  It should be easily recognisable.
\end{itemize}
\item Exposition repeat.  This is optional.  Traditionally, the
  exposition is repeated completely, the only possible difference
  being that the last few moments are varied to aid the smooth
  transition into the next section. (The `first' and `second time
  bars' which are in fact often several bars each.)  Variations are
  certainly possible.  Whether an exact repetition is a good idea is
  up to you.  Sometimes, only the first subject is repeated, and
  sometimes not in the exact same form. But note, however, that this
  repeat should not do the work of the next, development, section.
\item Development. This is the place where the conflict between the
  two themes is worked out in full.  In symphonies it is often very fiery
  and dramatic. Both themes are explored in detail in all possible keys
  (not only their home and away keys) sometimes against each other.
  The development section is usually the longest, typically half
  the length of the entire piece, possibly more.
\item Recapitulation. Traditionally, this is the same as the exposition
  except the second subject is now presented in the home key,
  not its away key, and transition material is altered appropriately.
  The important thing, however, is that both subjects return
  in a way similar to how they were introduced, but this time in the same key
  indicating that the conflict has been resolved.
\item (Optional) Coda.  A short sequence of material bringing the
  piece to an end.
\end{itemize}

\clearpage

\section{Suggestions}

You will want to think of objects as themes that can be varied
in the sense of having different keys (emotions, traits, characteristics).

Please make your two subjects recognisable as such, especially if your
first and second subjects are groups rather than individual objects.
It should be clear to the reader what any section of writing is
`about': the first subject, or the second subject, or the conflict
between the two (or possibly transitional material).

Exposition repeat: make the most of this if you can, and
don't feel shy about repeats in general.  There is
some wonderful use of a near-verbatim repeat near the beginning of
`A Portrait of the Artist as a Young Man' when Joyce recalls a
particularly horrible prank played on him by another schoolboy.
Repeats enable you to explore variations, both literary variations
and emotional variations, or (as in `Portrait') build up a sense
of emotion.  Note too that repeats are repeats of comparatively long
sections.  Repetitions of words and phrases within a section
can pall if over-used.

Resolution of the conflict:  take this in whatever form you
wish.  Indeed, there may be no resolution, or rather the resolution
that `two people agree to differ'.  The ending can be joyful, tragic,
funny, or whatever you want.

In later music, for example the first movement of Sibelius's first
symphony, the transition from the development to the recapitulation is
very smooth and `hidden' from the listener. Take from this that it is
not necessary that your piece should be divided into sections in the
way indicated in its final presentation, though that is how I suggest
you should be thinking of your piece as you write it and how you
organise your rough work.

By all means have a general plot or storyline in the back of your mind
but use sonata form structure to morph this into what is required.
The story most likely will be developed further as a result.
Or perhaps the sonata form structure is the starting point and suggests
to you a storyline.

I would imagine 25000--40000 words is about right for a prose piece, to
do this justice.  This need not be as daunting as it seems, since some
form of repetition-with-variation is encouraged.  For other media I
suggest 20-30 pages.  I am not setting any particular requirement on
length---only that I would imagine four or five of these pieces would
be enough for a single volume.

For all communications please email Sadie at \texttt{sadie@sadie.org.uk}.

Good luck!

\strut\hfill Sadie 

\end{document}
